%==============================================================================
\section{Thuật toán REINFORCE}
\label{sec:reinforce}
%==============================================================================

\begin{dinhnghia}[title={REINFORCE Algorithm}]
REINFORCE là thuật toán \textbf{policy gradient} trong RL, dựa trên phương pháp \textbf{Monte Carlo}.
Cách triển khai đơn giản: dùng \textbf{gradient ascent} để tăng trực tiếp phần thưởng tích lũy kỳ vọng bằng cách điều chỉnh tham số policy.
Thuật toán \textbf{không} cần mô hình môi trường $\Rightarrow$ thuộc nhóm \textbf{model-free}.
\end{dinhnghia}

\subsection{Khái niệm chính}

\begin{ynghia}[title={Policy Gradient Methods}]
REINFORCE thuộc lớp \textbf{policy gradient}: cải thiện policy bằng cách đi theo \textbf{gradient} của phần thưởng tích lũy kỳ vọng đối với tham số policy.
\end{ynghia}

\begin{ynghia}[title={Monte Carlo trong REINFORCE}]
REINFORCE là dạng Monte Carlo vì dùng \textbf{lấy mẫu (sampling)} để ước lượng các đại lượng mong muốn (return, gradient) từ \textbf{toàn bộ episode}, không bootstrap từng bước như TD.
\end{ynghia}

\subsection{Cách REINFORCE hoạt động}

\begin{dongco}[title={Lịch sử và target}]
Thuật toán do Ronald J. Williams giới thiệu (1992).
target: \textbf{tối đa hóa phần thưởng tích lũy kỳ vọng} bằng cách điều chỉnh tham số policy; agent học ra quyết định \textbf{tuần tự} trong môi trường.
\end{dongco}

\subsubsection{Lấy mẫu episode (Episode Sampling)}

\begin{phuongphap}[title={Bước 1 --- Lấy mẫu episode}]
Bắt đầu bằng cách lấy mẫu một \textbf{episode đầy đủ}: agent tương tác môi trường theo \textbf{policy hiện tại}.
Episode là chuỗi trạng thái, hành động và reward cho đến khi đạt \textbf{trạng thái kết thúc}.
\end{phuongphap}

\subsubsection{Quỹ đạo (Trajectory)}

\begin{phuongphap}[title={Bước 2 --- Ghi quỹ đạo}]
Agent ghi lại quỹ đạo tương tác:
\[
(s_1, a_1, r_1), (s_2, a_2, r_2), \ldots, (s_t, a_t, r_t)
\]
trong đó $s$ là trạng thái, $a$ là hành động, $r$ là reward tại mỗi bước.
\end{phuongphap}

\subsubsection{Return $G_t$}

\begin{phuongphap}[title={Bước 3 --- Tính return}]
Return $G_t$ là \textbf{phần thưởng tích lũy kỳ vọng} từ thời điểm $t$ trở đi:
\[
G_t = r_t + \gamma r_{t+1} + \gamma^2 r_{t+2} + \cdots
\]
với $\gamma$ là hệ số chiết khấu.
\end{phuongphap}

\subsubsection{Policy gradient}

\begin{phuongphap}[title={Bước 4 --- Gradient của return kỳ vọng}]
Tính gradient của return kỳ vọng theo tham số $\theta$ của policy $\pi_\theta$.
Theo nguyên lý policy gradient (Monte Carlo), cập nhật dựa trên \textbf{log-likelihood} của chuỗi hành động đã chọn:
\[
\nabla_\theta J(\theta) \approx \sum_{t=0}^{T-1} \nabla_\theta \log \pi_\theta(a_t \mid s_t)\, G_t
\]
\end{phuongphap}

\subsubsection{Cập nhật policy}

\begin{phuongphap}[title={Bước 5 --- Cập nhật tham số policy}]
Sau khi có gradient, cập nhật tham số theo hướng \textbf{tăng} reward kỳ vọng, ví dụ:
\[
\theta \leftarrow \theta + \alpha \sum_{t=0}^{T-1} \nabla_\theta \log \pi_\theta(a_t \mid s_t)\, G_t
\]
với $\alpha$ là learning rate.
\end{phuongphap}

\begin{luuy}[title={So với TD (Q-learning, SARSA)}]
Học TD (Q-learning, SARSA) thường nhấn mạnh reward \textbf{tức thời} và bootstrap.
REINFORCE cho agent học từ \textbf{toàn bộ chuỗi} trạng thái, hành động và reward của episode.
Lặp các bước trên cho đến khi episode (hoặc quá trình huấn luyện) kết thúc theo tiêu chí dừng.
\end{luuy}

\subsection{Ưu và nhược điểm}

\begin{tinhchat}[title={Ưu điểm REINFORCE}]
\begin{itemize}
  \item \textbf{Model-free}: phù hợp khi môi trường không biết hoặc khó mô hình hóa.
  \item \textbf{Đơn giản, trực quan}: dễ hiểu và triển khai.
  \item Xử lý được \textbf{không gian hành động chiều cao hoặc liên tục} --- khác nhiều phương pháp dựa trên giá trị rời rạc.
\end{itemize}
\end{tinhchat}

\begin{luuy}[title={Nhược điểm REINFORCE}]
\begin{itemize}
  \item \textbf{Phương sai cao}: ước lượng gradient biến động mạnh, học có thể chậm và \textbf{không ổn định}.
  \item \textbf{Dùng mẫu kém hiệu quả}: mỗi lần tính gradient thường cần bộ mẫu episode mới, kém tái sử dụng mẫu hơn các kỹ thuật dùng lại dữ liệu nhiều lần.
\end{itemize}
\end{luuy}
